%!TEX encoding = UTF-8 Unicode
%!TEX root = ../compendium.tex

\clearpage\null\thispagestyle{empty}
\vfill

{
\setlength{\parindent}{0pt}
\emph{Editor}: \ifswedish Björn Regnell \else Björn Regnell \fi \\

%  LIST OF CONTRIBUTORS to https://github.com/lunduniversity/introprog
%    Please contact bjorn.regnell@cs.lth.se if you think you should be
%    on this list, or make a pull request with an update of file briefly
%    describing your contribution in the commit text.
%    This work is licenced under CC-BY-SA-4.0.
%!TEX encoding = UTF-8 Unicode
%!TEX root = compendium/compendium.tex
\hyphenation{Borg-lund Da-ne-bjer Grampp Palm-qvist Ravn-borg Ro-sen-qvist Schrei-ter Wih-lan-der Nordahl Patrik}
\emph{Contributors} in alphabetical order:

{\raggedright
% \hyphenpenalty=10000
% \exhyphenpenalty=10000
% \vspace{0.5em}
% \noindent 

Alexander Sandström,
Anders Buhl,
André Philipsson Eriksson,
Anna Axelsson,
Anna Palmqvist Sjövall,
Annie Predel,
Anton Andersson,
Benjamin Lindberg,
Björn Regnell,
Casper Schreiter,
Cecilia Lindskog,
Dag Hemberg,
Elias Åradsson,
Elliot Bräck,
Elsa Cervetti Ogestad,
Emelie Engström,
Emil Wihlander,
Erik Bjäreholt,
Erik Grampp,
Evelyn Beck,
Felix Ohrgren,
Fredrik Danebjer,
Fritjof Bengtsson,
Gustav Cedersjö,
Henrik Olsson,
Hjalmar Rutberg,
Hussein Taher,
Jakob Hök,
Jakob Sinclair,
Johan Ekberg,
Johan Ravnborg,
Johannes Nydahl,
Jonas Danebjer,
Jos Rosenqvist,
Lovisa Löfgren, 
Maj Stenmark,
Maria Kulesh,
Mattias Nordahl,
Maximilian Waldenfeldt Uggla,
Melker Widén,
Måns Magnusson,
Nicholas Boyd Isacsson,
Niklas Sandén,
Oliver Levay,
Oliver Persson,
Oscar Sigurdsson,
Oskar Berg,
Oskar Widmark,
Patrik Persson,
Per Holm,
Philip Sadrian,
Sandra Nilsson,
Sebastian Hegardt,
Simon Persson,
Stefan Jonsson,
Theodor Lundqvist,
Tim Borglund,
Tom Postema,
Valthor Halldorsson,
Viggo Bergdahl,
Viktor Claesson,
Wilhelm Wanecek,
William Karlsson.
\par
}

~\\ \newline

\emph{Home}: \url{https://lunduniversity.github.io/pgk} \newline

\emph{Repo}: \url{https://github.com/lunduniversity/introprog} \\ \newline

\ifswedish Detta kompendium är under ständig utveckling. \else This compendium is ongoing work. \fi\\ \ifswedish \textbf{Bidrag välkomnas!} \else \textbf{Contributions are welcome!}\fi \\
\emph{Contact}: \url{bjorn.regnell@cs.lth.se}
\\ \newline

%\emph{Cover art}: Björn Regnell (inspired by Poul Ströyer's illustration of Lennart Hellsing's lyrics to  the childrens song ''Herr Gurka'' with music by Knut Brodin)\\ \newline

\emph{Versions:} \\
Scala \ScalaVersion \\
JDK \JDKVersion\\
\href{https://github.com/lunduniversity/introprog-scalalib/}{introprog-scalalib} \LibVersion \\ 

\vfill

You can use this work if you respect this \textbf{LICENSE}: CC BY-SA 4.0 \\
\url{http://creativecommons.org/licenses/by-sa/4.0/} \\
\ifswedish
\textbf{Dela \emph{inte} lösningar till laborationer och projekt.}
\else 
\textbf{Do \emph{not} distribute solutions to lab assignments and projects.}
\fi
\\ \newline
Copyright \copyright~ 2015-\CurrentYear. \\
Björn Regnell, Dept. of Computer Science, LTH, Lund University.\\
}
